\section{Nutzergruppen und Stakeholder}
\label{sec:03-04_user-groups-and-stakeholders}

Die Plattform wird von mehreren Nutzergruppen mit unterschiedlichen Rollen, Nutzungshäufigkeiten und Informationsbedarfen verwendet. Dazu zählen operative Mitarbeitende in Projekt- und Programmteams, koordinierende Rollen mit übergreifendem Informationsbedarf sowie Führungskräfte, die vor allem aggregierte, aktuelle und entscheidungsrelevante Informationen benötigen. Darüber hinaus sind Stakeholder aus IT, Organisationsentwicklung und interner Kommunikation beteiligt, die sowohl technische als auch strategische Anforderungen an die Plattform stellen, etwa in Bezug auf Systemintegration, Governance und die Unterstützung organisationaler Veränderungsprozesse \cite{308:DWGSegmentation,112:StepTwoPersonas}.

Die Analyse zeigte, dass diese Gruppen die Plattform auf sehr unterschiedliche Weise nutzen. Während einige Nutzende regelmäßig und tiefgehend mit Inhalten arbeiten, diese pflegen, strukturieren oder weiterentwickeln, greifen andere lediglich punktuell oder situativ auf Informationen zu, beispielsweise zur schnellen Recherche einzelner Dokumente, Statusinformationen oder Begriffe. Diese Unterschiede betreffen nicht nur die Nutzungshäufigkeit, sondern auch den Grad der inhaltlichen Verantwortung und Mitgestaltung innerhalb der Plattform.

Zur systematischen Einordnung dieser Heterogenität wurden die identifizierten Nutzergruppen entlang der Dimensionen Nutzungshäufigkeit der Plattform sowie Verantwortung für Inhalte analysiert und gegenübergestellt. Abbildung~\ref{fig:g-09_user-segmentation} visualisiert diese Differenzierung und verdeutlicht die Spannbreite zwischen primär konsumierenden und stark gestaltenden Nutzungsrollen sowie zwischen sporadischer und intensiver Nutzung.

\image{H}{\defaultImageSize}{g-09_user-segmentation}{
    Einordnung zentraler Nutzergruppen entlang der Dimensionen Nutzungshäufigkeit der Plattform und Verantwortung für Inhalte}

Die Darstellung macht sichtbar, dass Mitarbeitende innerhalb der Programmorganisation zu den intensiven Nutzenden zählen, die häufig auch inhaltlich gestaltend tätig sind, etwa durch das Erstellen, Aktualisieren oder Strukturieren von Inhalten. Mitarbeitende außerhalb der Programmorganisation nutzen die Plattform hingegen überwiegend gelegentlich und konsumierend, meist zur Informationsbeschaffung bei konkretem Bedarf. Führungskräfte nehmen eine Sonderrolle ein: Ihre Nutzung ist häufig nicht dauerhaft, jedoch mit hoher inhaltlicher und strategischer Verantwortung verbunden, da sie Inhalte steuern, freigeben oder als Referenz für Kommunikation und Entscheidungsprozesse heranziehen.

Diese Differenzierung wurde in der bestehenden Plattformstruktur nur unzureichend berücksichtigt. Insbesondere Gelegenheitsnutzende sowie neue Mitarbeitende trafen auf erhöhte Einstiegshürden, da Navigation, Begriffe und Einstiege primär an den Routinen erfahrener und intensiver Nutzergruppen ausgerichtet waren. Forschung zu Intranets und Digital Workplaces zeigt, dass eine fehlende Berücksichtigung unterschiedlicher Nutzungstypen häufig zu Nutzungslücken, reduzierter Akzeptanz und ineffizienter Informationssuche führt. Entsprechend wird empfohlen, explizit mit Nutzersegmenten und Employee Personas zu arbeiten, um Informationsarchitektur, Inhalteinstiege und Navigationskonzepte systematisch an unterschiedlichen Nutzungsmustern auszurichten \cite{113:NNGIntranetGuidelines,308:DWGSegmentation,112:StepTwoPersonas}.
